\documentclass[12pt]{article}
\usepackage{amsmath,amssymb,amsfonts}
\usepackage[margin=1in]{geometry}

\begin{document}

\textbf{Chapter 6, Problem 21.} \textit{Morera's theorem.} It states that if $f(z)$ is continuous in a region $R$ and if $\oint f(z)\,dz = 0$ for every closed contour in $R$, then $f(z)$ is analytic inside $R$. Prove this theorem.

\bigskip
\textbf{Solution.}

\textbf{Strategy:} We show that $f$ has an antiderivative, which is analytic, and then use that the derivative of an analytic function is analytic.

\textbf{Step 1: Construct an antiderivative.}

Since $\oint_C f(z)\,dz = 0$ for every closed contour in $R$, the integral of $f$ between any two points is path-independent. Fix a point $z_0 \in R$ and define:
\[
F(z) = \int_{z_0}^{z} f(\zeta)\,d\zeta,
\]
where the integral is along any path in $R$ from $z_0$ to $z$. This is well-defined by the path-independence.

\textbf{Step 2: Show $F'(z) = f(z)$.}

For $z, z+h \in R$:
\[
F(z+h) - F(z) = \int_z^{z+h} f(\zeta)\,d\zeta.
\]

We can take the straight-line path from $z$ to $z+h$. Then:
\[
\frac{F(z+h)-F(z)}{h} - f(z) = \frac{1}{h}\int_z^{z+h}[f(\zeta)-f(z)]\,d\zeta.
\]

Since $f$ is continuous, for any $\epsilon > 0$ there exists $\delta > 0$ such that $|f(\zeta)-f(z)| < \epsilon$ whenever $|\zeta - z| < \delta$. For $|h| < \delta$:
\[
\left|\frac{F(z+h)-F(z)}{h} - f(z)\right| \leq \frac{1}{|h|}\cdot\epsilon\cdot|h| = \epsilon.
\]

Therefore $F'(z) = f(z)$ exists for all $z \in R$.

\textbf{Step 3: Conclude $f$ is analytic.}

Since $F'(z) = f(z)$ exists everywhere in $R$, $F(z)$ is analytic in $R$. But a fundamental theorem of complex analysis states that the derivative of an analytic function is also analytic. Therefore $f(z) = F'(z)$ is analytic in $R$.

\end{document}
